% 图契约
% purpose: 读者应能回答——为何同一材料壳层的 \xi 不随收缩改变，以及两项为何逐点抵消
% topology: geometry
% reading_direction: coordinate-based
% node_roles: 横截面、材料壳层（两时刻）、固相位移、干基不变量
% edge_roles: 固相内移、半径对应、推导结论
% paper_context: 问题四动边界推导一节，紧随坐标映射示意，栏宽插入
% visual_encoding: 位置承担收缩，填充区分两时刻壳层，线型区分实体/收缩后边界，强调色仅标注结论
% style_family: G2
% annotation_policy: minimal
% result_policy: symbolic
% final_width: column
% print_mode: both
% MH-TIKZ-ABSOLUTE-OK: 坐标即柱体横截面的真实半径位置
\documentclass[border=3pt,11pt]{standalone}
\usepackage{ctex}
\usepackage{amsmath}
\usepackage{tikz}
\usetikzlibrary{arrows.meta,calc}

\definecolor{paperInk}{HTML}{111111}
\definecolor{paperAccent}{HTML}{3B549C}
\definecolor{paperGuide}{HTML}{4A4A4A}
\definecolor{paperPanel}{HTML}{E4E9F5}
\definecolor{paperCell}{HTML}{C6D0E8}

\begin{document}
\adjustbox{max width=\textwidth}{%
\begin{tikzpicture}[
  line cap=round,
  solid_body/.style={draw=paperInk, line width=0.9pt},
  thinline/.style={draw=paperInk, line width=0.55pt},
  guide/.style={draw=paperGuide, dashed, line width=0.5pt},
  emphasis/.style={draw=paperAccent, dash pattern=on 3pt off 1.6pt, line width=0.9pt},
  flux/.style={-{Stealth[length=4.5pt]}, draw=paperInk, line width=0.8pt},
  slbl/.style={text=paperInk, font=\footnotesize},
  acclbl/.style={text=paperAccent, font=\footnotesize}
]

% ================ 时刻 t 的材料壳层 \xi\in[0.6,0.8]，R(t)=3.0 ================
\fill[paperPanel] (2.4,0) arc[start angle=0, end angle=180, radius=2.4]
  -- (-1.8,0) arc[start angle=180, end angle=0, radius=1.8] -- cycle;
\draw[thinline] (2.4,0) arc[start angle=0, end angle=180, radius=2.4];
\draw[thinline] (1.8,0) arc[start angle=0, end angle=180, radius=1.8];

% ================ 时刻 t+\Delta t 的同一壳层，R 缩至 2.1 ================
\fill[paperCell] (1.68,0) arc[start angle=0, end angle=180, radius=1.68]
  -- (-1.26,0) arc[start angle=180, end angle=0, radius=1.26] -- cycle;
\draw[emphasis] (1.68,0) arc[start angle=0, end angle=180, radius=1.68];
\draw[emphasis] (1.26,0) arc[start angle=0, end angle=180, radius=1.26];

% ================ 外表面：两时刻 ================
\draw[solid_body] (3.0,0) arc[start angle=0, end angle=180, radius=3.0];
\draw[emphasis] (2.1,0) arc[start angle=0, end angle=180, radius=2.1];
\draw[thinline] (-3.4,0) -- (3.4,0);

% ================ 固相内移 ================
\draw[flux] (1.05,1.818) -- (0.735,1.273);
\draw[flux] (-1.05,1.818) -- (-0.735,1.273);

% ================ 两时刻壳层标记 ================
\node[slbl] at (1.818,1.05) {$t$};
\node[slbl] at (-1.273,0.735) {$t+\Delta t$};

% ================ 半径标注（左侧） ================
\draw[guide] (-3.0,0) -- (-3.0,-0.22);
\node[slbl] at (-3.0,-0.38) {$R(t)$};
\draw[guide] (-2.1,0) -- (-2.1,-0.88);
\node[slbl] at (-1.85,-1.05) {$R(t+\Delta t)$};

% ================ 物质坐标标注（右侧） ================
\draw[guide] (2.4,0) -- (2.4,-0.22);
\node[slbl] at (2.4,-0.38) {$\xi=0.8$};
\draw[guide] (1.8,0) -- (1.8,-0.88);
\node[slbl] at (1.8,-1.05) {$\xi=0.6$};

% ================ 侧向注解 ================
\node[slbl, align=center] at (-3.4,2.35) {固相内移\\$v_{s}<0$};
\node[slbl] at (3.55,2.35) {$\xi$ 不随收缩改变};

% ================ 干基不变量与抵消结论 ================
\node[slbl] at (0,3.52)
  {壳层干基质量 $\mathrm{d}m_{s}=\rho_{s}\,2\pi rL\,\mathrm{d}r$ 沿运动不变};

\node[acclbl, align=center] at (0,-2.05)
  {$\Rightarrow\ \rho_{s}R^{2}=$ 常数，\ $\xi=r/R(t)$ 为物质坐标\\
   网格漂移项与固相对流项逐点抵消};

\end{tikzpicture}
}%
\end{document}
